% PyMORGAN — Data processing and user manual.
% Build with:  pdflatex main.tex  (run twice, for the table of contents)
% Requires Python >= 3.12 and uv (https://docs.astral.sh/uv/)
\documentclass[11pt,a4paper]{article}

\usepackage[UKenglish]{isodate} % To fix the date
\usepackage[T1]{fontenc}
\usepackage{tgtermes}

\usepackage[utf8]{inputenc}
\usepackage[margin=2.5cm]{geometry}
\usepackage{amsmath,amssymb}
\usepackage{graphicx}
\usepackage{xcolor}
\usepackage{booktabs}
\usepackage{listings}
\usepackage[hidelinks]{hyperref}

\usepackage{microtype}        % best single improvement
	\emergencystretch=3em         % last-resort stretch to avoid overfull
	\tolerance=800                % default 200; higher = more flexible
	\hbadness=1000                % silence minor warnings

% --- Code listing style ----------------------------------------------------
\definecolor{codegray}{gray}{0.95}
\definecolor{codekw}{rgb}{0.10,0.30,0.65}
\definecolor{codestr}{rgb}{0.20,0.50,0.20}
\definecolor{codecom}{rgb}{0.45,0.45,0.45}
\lstdefinestyle{py}{
  language=Python,
  backgroundcolor=\color{codegray},
  basicstyle=\ttfamily\small,
  keywordstyle=\color{codekw}\bfseries,
  stringstyle=\color{codestr},
  commentstyle=\color{codecom}\itshape,
  showstringspaces=false,
  breaklines=true,
  frame=single,
  rulecolor=\color{codegray},
  columns=fullflexible,
}
\lstset{style=py}

% --- Convenience macros ----------------------------------------------------
\newcommand{\code}[1]{\texttt{#1}}
\newcommand{\dA}{\ensuremath{\Delta A}}

\title{\textbf{PyMORGAN}\\[0.3em]
  \large Multidimensional Optical spectroscopy Graphical Analysis iNterface\\
  \large Data processing and user manual}
\author{Dr. Ricardo J. Fernández-Terán}
\date{\today}

\begin{document}
\maketitle
\tableofcontents
\bigskip

\section{Introduction}
PyMORGAN provides a unified interface to load, process and plot ultrafast
time-resolved spectroscopy data. The package is organised by dimensionality:
the \code{oneD} module handles one-dimensional data (transient absorption /
pump--probe, time-resolved IR, fluorescence up-conversion), while \code{twoD}
targets two-dimensional data (2D-IR, 2D-ES, 2D-VE, 2D-EV), from raw
interferograms through phasing and Fourier transformation to the
spectral-diffusion and Kubo lineshape analyses. A common design
runs through both: a single dataset object wraps the raw arrays and exposes
processing and plotting methods, data formats are resolved through a loader
registry, and all presentation choices are centralised in a settings object.

This manual documents both the data-processing model and the programming
interface. Chapter~\ref{sec:installation} covers installation;
chapter~\ref{sec:oneD} covers the one-dimensional workflow;
chapter~\ref{sec:twoD} covers the two-dimensional workflow, including the
spectral-diffusion analysis; chapter~\ref{sec:steadyState} covers steady-state
absorption, emission and excitation spectra; chapter~\ref{sec:settings}
documents every setting and the configuration file;
chapter~\ref{sec:calibration} documents the spectrometer and pulse shaper
wavelength calibration engine; and chapter~\ref{sec:extending} (optional)
explains how to extend the GUI and add a parser for a new data format.

\section{Installation}\label{sec:installation}

PyMORGAN requires \textbf{Python~3.12 or later}. It can be installed directly from \href{https://pypi.org/project/pymorgan/}{PyPI} or installed from source using \href{https://docs.astral.sh/uv/}{\textbf{uv}}.

% ---------------------------------------------------------------------------
\subsection{Quick install via PyPI (Recommended)}\label{sec:install-pypi}

To install the latest official release into your active Python environment:

\begin{lstlisting}[language=bash, style=py]
pip install pymorgan
# or with uv:
uv pip install pymorgan

# Register bundled fonts into matplotlib:
pymorgan-install-fonts
\end{lstlisting}

\noindent You can also launch the graphical interface directly without prior installation using \code{uvx}:

\begin{lstlisting}[language=bash, style=py]
uvx --from pymorgan pymorgan-gui
\end{lstlisting}

% ---------------------------------------------------------------------------
\subsection{Installation from Source}\label{sec:install-source}

The steps below describe the full setup for development or source installation.

\paragraph{Step 1 — Install Python 3.12+.}

\paragraph{Windows.}
Download and run the official installer from
\url{https://www.python.org/downloads/}. Choose \textbf{Python 3.12} or
later, and tick \emph{Add Python to PATH} on the first installer page.

Alternatively, use the \textbf{Windows Store} or
\href{https://github.com/indygreg/python-build-standalone/releases}{Python
Standalone Builds} (recommended for reproducibility).

After the installer finishes, verify Python is on \code{PATH}:

\begin{lstlisting}[language=bash, style=py]
python --version        # should print Python 3.12.x or later
\end{lstlisting}

\paragraph{macOS.}
\begin{lstlisting}[language=bash, style=py]
brew install python@3.12          # Homebrew
# or download from https://www.python.org/downloads/
\end{lstlisting}

\paragraph{Linux (Debian/Ubuntu).}
\begin{lstlisting}[language=bash, style=py]
sudo apt update
sudo apt install python3.12 python3.12-venv python3.12-dev
\end{lstlisting}

% ---------------------------------------------------------------------------
\subsection{Step 2 — Install uv}\label{sec:install-uv}

uv is a single self-contained binary that replaces \code{pip},
\code{venv}, \code{pip-tools} and \code{virtualenv} in one tool. Install
it once, globally:

\paragraph{Windows (PowerShell).}
\begin{lstlisting}[language=bash, style=py]
powershell -ExecutionPolicy ByPass -c "irm https://astral.sh/uv/install.ps1 | iex"
\end{lstlisting}

\paragraph{macOS / Linux.}
\begin{lstlisting}[language=bash, style=py]
curl -LsSf https://astral.sh/uv/install.sh | sh
\end{lstlisting}

\paragraph{Via \code{pip} (any OS).}
\begin{lstlisting}[language=bash, style=py]
pip install uv
\end{lstlisting}

After installing, open a new terminal window and confirm:
\begin{lstlisting}[language=bash, style=py]
uv --version     # e.g. uv 0.6.x
\end{lstlisting}

% ---------------------------------------------------------------------------
\subsection{Step 3 — Obtain the PyMORGAN source}\label{sec:clone}

Clone the repository (or copy the folder to a location of your choice):

\begin{lstlisting}[language=bash, style=py]
git clone https://github.com/RJFernandezTeran/PyMORGAN.git
cd PyMORGAN
\end{lstlisting}

If you do not use git, download and extract the zip archive from GitHub,
then open a terminal inside the extracted folder.

% ---------------------------------------------------------------------------
\subsection{Step 4 — Create the virtual environment and install}\label{sec:venv}

From the repository root:

\begin{lstlisting}[language=bash, style=py]
uv venv                       # creates .venv/ with the system Python (>= 3.12)
uv pip install -e ".[dev]"    # editable install + ruff + pytest
uv run pymorgan-install-fonts # copy bundled fonts into matplotlib
\end{lstlisting}

\noindent\textbf{What each line does:}
\begin{itemize}
  \item \code{uv venv} — creates a \code{.venv/} virtual environment in the
    current directory.  Pass \code{--python 3.12} to pin a specific
    interpreter if multiple versions are installed.
  \item \code{uv pip install -e ".[dev]"} — installs all core dependencies
    (Table~\ref{tab:dependencies}), plus the development extras
    (\code{ruff} linter and \code{pytest} test runner), in \emph{editable}
    mode. Editable mode means changes to the source files take effect
    immediately without reinstalling.
  \item \code{uv run pymorgan-install-fonts} — copies the bundled
    \code{TgTermes} and other fonts used by the plot style profiles into
    matplotlib's font cache. Run this once after a fresh install.
\end{itemize}

\noindent If you only need to run the GUI (no development work):

\begin{lstlisting}[language=bash, style=py]
uv pip install -e .           # core only (no ruff / pytest)
uv run pymorgan-install-fonts
\end{lstlisting}

% ---------------------------------------------------------------------------
\subsection{Step 5 — Verify the installation}\label{sec:verify}

Confirm the package imports and reports its version:

\begin{lstlisting}[language=bash, style=py]
uv run python -c "import pymorgan; print(pymorgan.__version__)"
\end{lstlisting}

Launch the graphical interface:

\begin{lstlisting}[language=bash, style=py]
uv run pymorgan-gui
\end{lstlisting}

Run the full test suite (requires the \code{dev} extras):

\begin{lstlisting}[language=bash, style=py]
uv run pytest        # 112 tests, all should pass
\end{lstlisting}

Run the end-to-end smoke check:

\begin{lstlisting}[language=bash, style=py]
uv run python scripts/run_checks.py
\end{lstlisting}

% ---------------------------------------------------------------------------
\subsection{Dependencies}\label{sec:dependencies}

Table~\ref{tab:dependencies} lists the core runtime dependencies. All are
installed automatically by \code{uv pip install -e .}; no manual
\code{pip install} is needed.

\begin{table}[h]
  \centering
  \begin{tabular}{lll}
    \toprule
    Package & Minimum version & Purpose \\
    \midrule
    \code{numpy}           & 2.5.2  & Array operations \\
    \code{scipy}           & 1.18.0 & Signal processing and fitting \\
    \code{matplotlib}      & 3.11.1 & Plotting back-end \\
    \code{pandas}          & 3.0.5  & Tabular data and I/O \\
    \code{h5py}            & 3.16.0 & HDF5 file I/O (some loaders) \\
    \code{mpl-axes-aligner}& 1.3    & Aligned multi-axis figures \\
    \code{tomlkit}         & 0.15.1 & Settings file read/write \\
    \code{PyQt6}           & 6.11.0 & Graphical interface framework \\
    \code{PyQt6\_sip}      & 13.12.0& Qt bindings interface \\
    \bottomrule
  \end{tabular}
  \caption{Core runtime dependencies (minimum version floors as of August 2026).
    Exact versions are pinned in \code{uv.lock}.}
  \label{tab:dependencies}
\end{table}

\begin{table}[h]
  \centering
  \begin{tabular}{lll}
    \toprule
    Package & Minimum version & Purpose \\
    \midrule
    \code{ruff}      & 0.16.2 & Linter and formatter \\
    \code{pytest}    & 9.1.1  & Test runner \\
    \code{pytest-qt} & 4.5.0  & Qt widget testing helpers \\
    \bottomrule
  \end{tabular}
  \caption{Development extras (\code{uv pip install -e ".[dev]"}).}
  \label{tab:dev-dependencies}
\end{table}

% ---------------------------------------------------------------------------
\subsection{Keeping packages up to date}\label{sec:upgrade}

To upgrade all dependencies to the latest compatible versions:

\begin{lstlisting}[language=bash, style=py]
uv lock --upgrade     # resolve latest versions, update uv.lock
uv sync               # install the resolved versions into .venv
\end{lstlisting}

% ---------------------------------------------------------------------------
\subsection{Console scripts}\label{sec:scripts}

Four console scripts are installed with the package:

\begin{table}[h]
  \centering
  \begin{tabular}{ll}
    \toprule
    Script & Purpose \\
    \midrule
    \code{pymorgan-gui}           & Launch the PyQt6 graphical interface \\
    \code{pymorgan-install-fonts} & Install bundled fonts into matplotlib \\
    \code{pymorgan-edit-gui}      & Open the GUI layout in Qt Designer \\
    \code{pymorgan-settings}      & Standalone settings editor (no data loaded) \\
    \bottomrule
  \end{tabular}
  \caption{Console scripts installed with the package.}
  \label{tab:scripts}
\end{table}

All scripts are run via \code{uv run <script>} (e.g.\
\code{uv run pymorgan-gui}), or directly if the environment is activated.

% ---------------------------------------------------------------------------
\subsection{GUI layout and compiled UI}\label{sec:compiled-ui}

The GUI layout is defined in
\code{src/pymorgan/gui/main\_window.ui} (Qt Designer XML). For performance,
PyMORGAN also ships a pre-compiled Python translation,
\code{main\_window\_ui.py}, generated by \code{pyuic6}. Loading the
compiled module is approximately $1.5\times$ faster than parsing the XML at
runtime.

\textbf{No manual action is needed when editing the \code{.ui} file.} When
\code{main\_window.ui} is saved in Qt Designer and the GUI is next launched,
PyMORGAN detects that the compiled module is stale (via file modification
times) and regenerates it automatically using \code{pyuic6} before loading
the window. If \code{pyuic6} is unavailable or fails, the raw XML is parsed
as a transparent fallback.

To force a manual recompile (e.g.\ for CI):

\begin{lstlisting}[language=bash, style=py]
uv run pyuic6 src/pymorgan/gui/main_window.ui \
              -o src/pymorgan/gui/main_window_ui.py
\end{lstlisting}

% ---------------------------------------------------------------------------
\subsection{Activating the environment (optional)}\label{sec:activate}

All commands in this manual use the \code{uv run} prefix, which
automatically targets the project's \code{.venv} without requiring
activation. If you prefer an activated shell (e.g.\ for interactive work
in a terminal):

\begin{lstlisting}[language=bash, style=py]
# Linux / macOS
source .venv/bin/activate

# Windows (PowerShell)
.venv\Scripts\Activate.ps1

# Windows (Command Prompt)
.venv\Scripts\activate.bat
\end{lstlisting}

Once activated, commands can be run without the \code{uv run} prefix
(e.g.\ \code{pymorgan-gui} instead of \code{uv run pymorgan-gui}).

\section{One-dimensional data}\label{sec:oneD}

\subsection{Pipeline overview}
The one-dimensional branch is organised as a \emph{load $\to$ process $\to$
plot $\to$ analyse} pipeline. Each stage is a module under
\code{pymorgan.oneD}, and the \code{Dataset1D} object ties them together:
it holds the raw arrays and exposes one set of methods per stage. The stage
functions take the dataset as their first argument, so they remain reusable on
their own and new routines integrate with the whole pipeline. Table~\ref{tab:stages}
maps each stage to its module and entry points.

\begin{table}[h]
  \centering
  \begin{tabular}{lll}
    \toprule
    Stage & Module & Entry point(s) \\
    \midrule
    load    & \code{pymorgan.oneD.load}    & \code{pm.load\_1D}, \code{Dataset1D.from\_file} \\
    process & \code{pymorgan.oneD.process} / \code{chirp} & \code{Dataset1D.background\_correct}, \code{apply\_chirp\_correction} \\
    plot    & \code{pymorgan.oneD.plot}    & \code{Dataset1D.plot\_*} \\
    analyse & \textit{(PyRATE-TA)} & \textit{separate project} \\
    \bottomrule
  \end{tabular}
  \caption{The 1-D pipeline stages and where they live.}
  \label{tab:stages}
\end{table}

\subsection{Data model}
A one-dimensional measurement is represented by a \code{Dataset1D} object,
which bundles the arrays produced by a loader. The signal is stored as a
three-dimensional array to generalise over multiple detectors,
\begin{equation}
  Z \in \mathbb{R}^{N_\tau \times N_\lambda \times N_d},
\end{equation}
indexed by delay $\tau$, probe coordinate $\lambda$ (wavelength, wavenumber or
energy) and detector $d$. The accompanying axes are the delay vector
$\tau \in \mathbb{R}^{N_\tau}$ and the probe vector
$\lambda \in \mathbb{R}^{N_\lambda}$. Table~\ref{tab:fields} lists the seven
fields every loader returns.

\begin{table}[h]
  \centering
  \begin{tabular}{lll}
    \toprule
    Field & Symbol / shape & Meaning \\
    \midrule
    \code{Zavg\_R} & $N_\tau \times N_\lambda \times N_d$ & averaged signal \\
    \code{delays}  & $N_\tau$ & delay axis \\
    \code{probe}   & $N_\lambda$ & probe (spectral) axis \\
    \code{Units}   & dict & unit labels (see \S\ref{sec:settings}) \\
    \code{Nscans}  & scalar & number of scans (\code{nan} if unknown) \\
    \code{Zss\_R}  & $N_\tau \times N_\lambda \times N_d \times N_s$ & single-scan signal \\
    \code{Zstdv}   & $N_\tau \times N_\lambda$ & standard deviation \\
    \bottomrule
  \end{tabular}
  \caption{Fields returned by a 1-D loader and stored on \code{Dataset1D}.}
  \label{tab:fields}
\end{table}

The most-processed signal is exposed through the \code{Z} property, which
returns the background-corrected array once a correction has been applied and
falls back to the raw array otherwise.

\subsection{Loading and the loader registry}
Data formats are decoupled from the dataset through a \emph{loader registry}
(\code{pymorgan.oneD.load}). A loader is any callable that maps a file path to
the seven fields of Table~\ref{tab:fields}. Loaders are registered under one or
more data-type names and resolved at load time:

\begin{lstlisting}
import pymorgan as pm

data = pm.load_1D("scan.pdat", data_type="PDAT")
print(pm.available_loaders())
\end{lstlisting}

The currently registered names are \code{PDAT} (implemented) together with the
instrument-specific formats \code{UniGE\_fsTA}, \code{UniGE\_nsTA},
\code{MESS\_TRIR}, \code{MESS\_TRUVIS}, \code{UniGE\_FLUPSold},
\code{UniGE\_FLUPSnew}, \code{Helios\_TA} and \code{UoS\_IRpp}, which are
recognised but raise \code{NotImplementedError} until their on-disk formats are
defined. A new format is added by decorating a reader:

\begin{lstlisting}
@pm.register_loader("MyInstrument")
def read_my_instrument(path):
    ...
    return Zavg_R, delays, probe, Units, Nscans, Zss_R, Zstdv
\end{lstlisting}

\paragraph{The PDAT format.}
A \code{.pdat} file is comma-separated. The first line carries the time and
probe units separated by an asterisk (e.g.\ \code{ps*cm\textasciicircum\{-1\}}).
The first data row holds the probe axis (with a leading $0$), and each
subsequent row begins with a delay value followed by the signal at every probe
coordinate. PDAT files store averaged data already in m\,OD, so no rescaling is
applied on read.

\subsection{Background correction}
The pre-time-zero background is estimated by averaging the signal over a delay
window $[\tau_\mathrm{min},\tau_\mathrm{max}]$ and subtracting it from every
delay,
\begin{equation}
  Z_C(\tau,\lambda,d) = Z(\tau,\lambda,d)
    - \big\langle Z(\tau',\lambda,d) \big\rangle_{\tau' \in [\tau_\mathrm{min},\tau_\mathrm{max}]},
\end{equation}
where the average is taken with \code{nan}-aware statistics
(\code{pymorgan.oneD.process}). The operation is applied through

\begin{lstlisting}
data.background_correct(tmin=-20, tmax=-5)          # apply
data.background_correct(tmin=-20, tmax=-5, do_correct=False)  # passthrough
\end{lstlisting}

With \code{do\_correct=False} the raw signal is copied unchanged into the
corrected fields, so downstream code can always read \code{Z} uniformly. The
method returns \code{self}, allowing calls to be chained.

\subsection{Solvent subtraction}\label{sec:solvent}
For measurements in solution, the pure-solvent response often has to be
subtracted. The module exposes two entry points.

\paragraph{Manual subtraction.} A scaled, time-shifted solvent trace is removed
from the sample signal:
\begin{lstlisting}
data.subtract_solvent(solvent, dt=0.0, scale=1.0)
\end{lstlisting}
\code{dt} shifts the solvent delay axis before interpolation (accounts for a
timing offset between the two measurements) and \code{scale} multiplies the
solvent signal before subtraction.

\paragraph{Automatic solvent fitting.} \code{fit\_solvent\_auto} models the
solvent contribution pixel by pixel using the solvent trace convolved with the
Instrument Response Function (IRF) and removes it:
\begin{lstlisting}
fit = data.fit_solvent_auto(solvent)
data.subtract_solvent_fit(fit)
\end{lstlisting}
Optional per-pixel fits for the time-zero offset and the amplitude scale can be
enabled through the \code{per\_pixel\_dt} and \code{per\_pixel\_scale} settings
(Section~\ref{sec:settings}); the defaults are \code{false} and \code{true},
respectively.

\subsection{Shockwave subtraction}\label{sec:shockwave}
In transient IR / UV-Vis measurements using flowing samples, acoustic shockwaves
generated by the pump beam can appear as a time-dependent artefact that is
correlated across \emph{all} probe wavelengths. The artefact is removed by
averaging a selected subset of detector pixels (or a probe-axis range) into a
one-dimensional shockwave trace and subtracting it from every pixel:
\begin{lstlisting}
from pymorgan.oneD.process import subtract_shockwave

Z_corrected = subtract_shockwave(Z, pixels=[0, 1, 2])   # pixel indices (0-based)
\end{lstlisting}

The GUI exposes this through the \textbf{Subtract Shockwave\ldots} button in the
1-D tab, which opens the \code{ShockwaveSubtractionDialog}. The dialog lets the
user choose between pixel-index and probe-axis-value selection, displays a live
preview of the proposed shockwave trace, and applies the correction on
confirmation (see \S\ref{sec:oneD:gui:dialogs}).

\subsection{Chirp (dispersion) correction}\label{sec:chirp}
Ultrafast time-resolved spectroscopy measurements often suffer from Group Delay Dispersion (GDD), commonly referred to as ``chirp''. As the probe pulse passes through various dispersive optical elements (such as white-light generation crystals, lenses, sample cell windows, and the solvent itself), different spectral components are temporally stretched. Consequently, the effective time-zero $t_0(\lambda)$ varies across the probe wavelength range $\lambda$.

PyMORGAN provides tools to characterise and correct this dispersion. The time-zero profile $t_0(\lambda)$ is modelled using a generalised Cauchy dispersion equation:
\begin{equation}
  t_0(\lambda) = c_0 + \sum_{n=1}^{N_C-1} c_n \left(\frac{\lambda_{\mathrm{ref}}}{\lambda}\right)^{2n},
\end{equation}
where $c_0, \dots, c_{N_C-1}$ are the dispersion coefficients, $\lambda_{\mathrm{ref}}$ is a reference wavelength (which defaults to the mean of the fitted range to make the coefficients dimensionless and of comparable magnitude), and $N_C$ is the number of terms (usually $6$).

To determine $t_0(\lambda)$ across the probe range, PyMORGAN implements three fit modes in the \code{pymorgan.oneD.chirp} module:
\begin{description}
  \item[Automatic] Fits the coherent artefact and early dynamics at every pixel independently. The model consists of a Gaussian Instrument Response Function (IRF), its optional first and second derivatives, a constant offset, and an optional erf-broadened exponential decay. PyMORGAN uses Variable Projection (VARPRO) to eliminate the linear amplitudes in closed form, optimizing only the nonlinear parameters $(t_0, \mathrm{FWHM}, \tau_{\mathrm{exp}})$ per pixel. A second pass uses smoothed parameters from the first pass as initial guesses to improve stability.
  \item[Step function] A simplified alternative that jointly fits a Gaussian IRF (+ derivatives) and a Heaviside step convolved with the same Gaussian over a narrow early-time window. This is robust when the early-time exponential decay is not needed or is ill-conditioned.
  \item[Manual] Fits the Cauchy dispersion equation directly to a set of user-picked $(\lambda, t_0)$ points (e.g., clicked interactively on the contour map in the GUI), bypassing the per-pixel IRF fitting.
\end{description}

Once $t_0(\lambda)$ is obtained, the Cauchy dispersion equation is fit using bound-constrained, iteratively reweighted linear least-squares (IRLS with Tukey's bisquare weights) to handle outliers and reject bad pixel fits.

To apply the correction, the dataset's delay axes are interpolated onto a common chirp-free delay grid. Two boundary handling strategies are available:
\begin{itemize}
  \item \code{crop}: Slices the delay axis to keep only the delay region where all wavelength channels are defined, discarding the out-of-bounds regions that become \code{NaN} after shifting.
  \item \code{extrapolate}: Fills out-of-bound values with the nearest valid boundary endpoint value, keeping the original delay range.
\end{itemize}

\begin{lstlisting}
# Fit chirp automatically and apply the correction
fit = data.fit_chirp_automatic(deriv_mask=(True, True), use_exp=True)
data.plot_chirp_diagnostics(fit)
data.apply_chirp_correction(fit, boundary_handling="crop")
\end{lstlisting}

\subsection{Plotting}
Five plotters live in \code{pymorgan.oneD.plot}; each is exposed both as a
\code{Dataset1D} method and as a stage-level function that takes the dataset as
its first argument, so a new view only needs \code{(data, ax, \dots)} and
inherits the whole pipeline's context. Each returns the matplotlib axis (the
kinetic view also returns the extracted traces). They are listed in
Table~\ref{tab:plotters}.

\begin{table}[h]
  \centering
  \begin{tabular}{ll}
    \toprule
    Method / function & View \\
    \midrule
    \code{plot\_contour}        & delay-vs-probe contour map \\
    \code{plot\_surface}        & filled delay-vs-probe surface \\
    \code{plot\_spectra}        & transient spectra at selected delays \\
    \code{plot\_kinetics}       & kinetic traces at selected probe positions \\
    \code{plot\_species\_spectra} & EAS/SAS from a global-analysis fit \\
    \bottomrule
  \end{tabular}
  \caption{The 1-D plotters (\code{pymorgan.oneD.plot}).}
  \label{tab:plotters}
\end{table}

\begin{lstlisting}
data.plot_contour(Zscale=20)                          # method form
data.plot_spectra([0.5, 1, 5, 20, 100], doSmooth=1)
ax, t, Y = data.plot_kinetics([2132, 2218], plotStyle="-")

from pymorgan.oneD import plot                          # stage-function form
plot.plot_contour(data, label_style="/", cmap_ID="Rd/Wh/Bu v2")
\end{lstlisting}

The label style, colourmap, time-axis scale and \dA{} unit convention are taken
from the active settings (\S\ref{sec:settings}); any of them can be overridden
per figure through the corresponding keyword argument.

For datasets with high dynamic range where weak features would otherwise be washed out by dominant bleach or ESA peaks, \code{plot_contour} provides non-linear $\operatorname{arcsinh}$ colour scaling via \code{Asinh=True} (also toggled in the GUI plot controls). Rather than modifying the underlying signal values, PyMORGAN applies an inverse hyperbolic sine colour normalization (\code{matplotlib.colors.AsinhNorm}) paired with non-linear $\sinh$-spaced contour levels:
\begin{equation}
  z_k = s \cdot \sinh(u_k), \quad u_k \in \left[-\operatorname{arcsinh}\left(\frac{|Z_{\min}|}{s}\right), \operatorname{arcsinh}\left(\frac{Z_{\max}}{s}\right)\right],
\end{equation}
where $s = \frac{p}{100} \cdot \max(|Z_{\min}|, |Z_{\max}|)$ is the linear crossover threshold set by \code{asinh_pct} ($p$, default $5.0\%$). Signals within $\pm s$ behave linearly, while larger amplitudes are compressed logarithmically. Because the underlying data array is not mutated, coordinates and hover values retain their physical units ($\mathrm{mOD}$), and colorbar ticks remain linear and clean (centred at zero and rounded to sensible numbers).

\subsection{Composing figures and displaying}
By default every plotter creates its own figure and axis. Passing an existing
axis through \code{ax=} draws into it instead, which is how panels are combined
into a layout or several datasets are overlaid on one axis. Repeated elements
(axis labels and, for the maps, the colorbar) can be switched off so they are
not drawn more than once; the toggles are listed in Table~\ref{tab:toggles}.

\begin{table}[h]
  \centering
  \begin{tabular}{lll}
    \toprule
    Keyword & Default & Plotters \\
    \midrule
    \code{show\_xlabel}        & \code{True} & all \\
    \code{show\_ylabel}        & \code{True} & all \\
    \code{show\_colorbar}      & \code{True} & \code{plot\_contour}, \code{plot\_surface} \\
    \code{show\_colorbar\_label} & \code{True} & \code{plot\_contour}, \code{plot\_surface} \\
    \bottomrule
  \end{tabular}
  \caption{Per-call toggles for composing multi-panel figures.}
  \label{tab:toggles}
\end{table}

\begin{lstlisting}
import matplotlib.pyplot as plt
fig, (axL, axR) = plt.subplots(1, 2)

# two maps side by side; suppress the duplicated colorbar / y-axis label
dataA.plot_contour(ax=axL, show_colorbar=False)
dataB.plot_contour(ax=axR, show_ylabel=False)

# overlay two datasets' transient spectra on a single axis
dataA.plot_spectra([1, 5, 20], ax=axL)
dataB.plot_spectra([1, 5, 20], ax=axL, show_ylabel=False)

pm.show_plots()   # wrapper around plt.show(); pm.show is an alias
\end{lstlisting}

When a plotter is given an existing axis it does not resize the parent figure,
so dropping a plot into a layout leaves the surrounding panels untouched.
\code{pm.show\_plots(block=False)} returns immediately (e.g.\ inside a GUI event
loop), and \code{pm.close\_plots()} closes all open figures.

\subsection{Kinetic analysis}\label{sec:analyse}
Kinetic analysis is \emph{not} part of PyMORGAN. Trace extraction and fitting
--- multi-exponential decays, and full global/target analysis with rate-matrix
models (the source of the EAS/SAS consumed by \code{plot\_species\_spectra})
--- live in the companion project \textbf{PyRATE-TA}, which shares this
loader/dataset layer.

PyMORGAN keeps the two ends of that seam. Traces are obtained from
\code{plot\_kinetics}, whose third return value is the extracted data itself:

\begin{lstlisting}
ax, t, Y = data.plot_kinetics([1950, 2000], plotStyle="-")
Y.shape                                           # [Ndelays x Ncuts]
\end{lstlisting}

and fitted spectra are rendered by \code{plot\_species\_spectra}, which takes
plain arrays (\code{Sfit}, \code{Taus}, \code{TauErr}, \code{isFixTau},
\code{modelType}) and is therefore independent of how the fit was produced.

The third piece of the seam is the noise. \code{Dataset1D.noise\_array()}
returns the per-point standard deviation --- from \code{Zstdv} when a sibling
\code{.pdatn} file was loaded, otherwise the spread across single scans, and
\code{None} when the dataset carries neither:

\begin{lstlisting}
sigma = data.noise_array()      # [Ndelays x Npixels x Ndetectors], or None
\end{lstlisting}

A weighted fit uses it as \code{1/sigma}, which is what turns a fit statistic
into a true reduced $\chi^2$; without it the statistic is a plain sum of
squared residuals. PyMORGAN owns the noise because it owns the loaders; PyRATE-TA
never estimates it.

\subsection{Settings and aesthetics}\label{sec:settings}
Presentation choices that are not intrinsic to the data are collected in a
single \code{Settings} object. It is framework-agnostic: the GUI binds to the
fields it describes, but the package never depends on the GUI. The fields are
summarised in Table~\ref{tab:settings}.

\begin{table}[h]
  \centering
  \begin{tabular}{lll}
    \toprule
    Field & Values & Effect \\
    \midrule
    \code{profile} & regular / poster / inset & selects a \code{.mplstyle} file \\
    \code{font\_scale} & float & multiplies the profile font sizes \\
    \code{label\_style} & \code{()} \code{[]} \code{/} (empty) & axis-label delimiter \\
    \code{delta\_a\_units} & \code{mOD} / \code{x1E3} & \dA{} label convention \\
    \code{cmap} & string & default contour colourmap \\
    \code{time\_axis\_scale} & symlog / log / lin & default time-axis scale \\
    \code{freq\_label} & Omega\_n / Omega\_PP / \dots & 2-D frequency-axis convention \\
    \code{x\_axis\_unit} & native / nm / cm-1 / eV / THz & spectral X-axis display unit \\
    \code{secondary\_axis} & bool & complementary X axis (top) \\
    \code{chirp\_boundary\_handling} & crop / extrapolate & boundary handling for chirp correction \\
    \bottomrule
  \end{tabular}
  \caption{Fields of the \code{Settings} object.}
  \label{tab:settings}
\end{table}

The \code{label\_style} field controls the delimiter so that an axis reads, for
example, ``Wavelength (nm)'', ``Wavelength [nm]'' or ``Wavelength / nm''. The
\code{delta\_a\_units} field switches the signal-axis label between
``\dA{} / m\,OD'' and ``\dA{} $\times 10^{3}$''; because $1\,\mathrm{mOD}
\equiv \dA \times 10^{3}$, this is a labelling choice only and never rescales
the stored data. The active \code{label\_style} delimiter is honoured for the
m\,OD label on every axis, including the contour colorbar (so ``/'' yields
``\dA{} / m\,OD''); the $\times 10^{3}$ form is always shown without a delimiter.

The style profiles wrap the existing hand-tuned \code{.mplstyle} files rather
than replacing them; \code{font\_scale} is applied as a light multiplicative
overlay on top. Text is always rendered with matplotlib's mathtext engine; the
LaTeX backend (\code{text.usetex}) is never enabled.

The \code{x\_axis\_unit} field converts the spectral (probe) X axis of the transient-spectra and contour plots between wavelength (nm), wavenumber (cm$^{-1}$), energy (eV) and frequency (THz); \code{native} keeps each dataset's stored unit. The conversion is physical (the axis values are recomputed, pivoting through the wavelength), not a relabelling. When \code{secondary\_axis} is enabled a second X axis is drawn on top in the \emph{complementary} unit: cm$^{-1}\leftrightarrow$nm, with nm used as the complement for eV and THz. Both fields default to the active settings and can be overridden per call, e.g. \code{data.plot\_spectra([1, 5, 20], x\_axis\_unit="nm", secondary\_axis=True)}. The earlier \code{dualScale} keyword is retained as a deprecated alias (a truthy value maps to \code{secondary\_axis=True}). Whenever a wavenumber axis (main or secondary) exceeds $5000\,\mathrm{cm^{-1}}$ its tick labels are divided by $1000$ and the axis is relabelled ``Wavenumber ($\times 10^{3}$ cm$^{-1}$)'', so UV/Vis wavenumbers stay legible; IR ranges are left unscaled. The axis data, stored limits and GUI spin boxes remain in true cm$^{-1}$ -- only the displayed tick labels and unit are scaled.

The \code{chirp\_boundary\_handling} field controls how out-of-bounds delays are treated during the chirp correction interpolation step, selecting between \code{crop} and \code{extrapolate} (see \S\ref{sec:chirp}).

Settings are read from and written to a TOML file, with comments preserved on
write so the file remains hand-editable after the GUI saves it:

\begin{lstlisting}
pm.load_settings("settings.toml")   # make the file's settings active
pm.update_settings(profile="poster", label_style="/")
pm.apply_style()                    # apply the profile to matplotlib
pm.save_settings("settings.toml")   # persist (keeps comments)
\end{lstlisting}

A module-level \emph{active} settings object is read by the plotting methods.
Temporary overrides can be scoped with a context manager:

\begin{lstlisting}
with pm.use_settings(delta_a_units="mOD"):
    data.plot_spectra([1, 5, 20])   # this figure only
\end{lstlisting}

\subsection{Worked example}
A complete script reproducing the standard workflow is provided in
\code{examples/pump\_probe\_1D.py}. In outline, the four stages read:

\begin{lstlisting}
import pymorgan as pm

pm.load_settings("settings.toml")
pm.apply_style()

data = pm.load_1D("testData/testTRIR.pdat", data_type="PDAT")  # load
data.background_correct(tmin=-20, tmax=-5)                       # process
data.plot_contour(Zscale=20, Yscale="lin")                      # plot
ax, t, Y = data.plot_kinetics([1980, 2030], plotStyle="-")      # plot + extract
\end{lstlisting}

\subsection{Graphical interface — 1-D dialogs}\label{sec:oneD:gui:dialogs}

Several interactive dialogs supplement the main 1-D tab.

\subsubsection{Chirp correction dialog}

The \textbf{Chirp\ldots} button opens \code{chirp\_dialog.py}, which guides
the user through chirp correction in three sub-panels:
\begin{enumerate}
  \item \textbf{Mode selection} --- choose Automatic, Step-function or Manual.
  \item \textbf{Parameter input} --- set the delay window, IRF FWHM, and
        (for Automatic mode) which derivative terms to include.
  \item \textbf{Diagnostics} --- the fitted $t_0(\lambda)$ curve overlaid on
        the contour map; the boundary-handling strategy (crop / extrapolate)
        is chosen here before the correction is applied.
\end{enumerate}
A progress dialog (\code{chirp\_progress\_dialog.py}) reports per-pixel fit
progress for the Automatic mode, which can be parallelised via the
\code{parallel\_fitting} setting (thread pool, process pool, or sequential).

\subsubsection{Solvent subtraction dialog}

The \textbf{Subtract Solvent\ldots} button opens a dialog in
\code{dialogs.py} that loads a pure-solvent scan, previews the subtraction,
and applies \code{fit\_solvent\_auto} or the manual \code{subtract\_solvent}
path according to the chosen settings.

\subsubsection{Shockwave subtraction dialog}

The \textbf{Subtract Shockwave\ldots} button opens
\code{shockwave\_dialog.py}. The user selects a set of pixels (by index or
probe value) whose average forms the shockwave reference trace, inspects a
live preview, and confirms. The correction calls \code{subtract\_shockwave}
and updates the dataset in place.

\section{Two-dimensional data}\label{sec:twoD}

The two-dimensional branch follows the same \emph{load $\to$ process $\to$
plot $\to$ analyse} pipeline as the one-dimensional one (chapter~\ref{sec:oneD}),
built around a single \code{Dataset2D} object. Table~\ref{tab:stages2d}
maps each stage to its module.

\begin{table}[h]
  \centering
  \begin{tabular}{lll}
    \toprule
    Stage & Module & Entry point(s) \\
    \midrule
    load    & \code{pymorgan.twoD.load}    & \code{pm.load\_2D}, \code{Dataset2D.from\_file} \\
    process & \code{pymorgan.twoD.process} & \code{Dataset2D.background\_correct} \\
    plot    & \code{pymorgan.twoD.plot}    & \code{Dataset2D.plot\_map} \\
    analyse & \code{pymorgan.twoD.analyse} & \code{Dataset2D.slice\_at}, \code{diagonal} \\
    \bottomrule
  \end{tabular}
  \caption{The 2-D pipeline stages and where they live.}
  \label{tab:stages2d}
\end{table}

\subsection{Data model}
A two-dimensional measurement is a stack of 2-D frequency--frequency maps, one
per population time $t_2$. The signal is stored as a cube
\begin{equation}
  Z \in \mathbb{R}^{N_{\text{pump}} \times N_{\text{probe}} \times N_{t_2}},
\end{equation}
indexed by the pump frequency $\omega_1$, the probe frequency $\omega_3$ and the
population time $t_2$. The fields held by \code{Dataset2D} are listed in
Table~\ref{tab:fields2d}.

\begin{table}[h]
  \centering
  \begin{tabular}{lll}
    \toprule
    Field & Shape & Meaning \\
    \midrule
    \code{Z}      & $N_{\text{pump}} \times N_{\text{probe}} \times N_{t_2}$ & signal cube \\
    \code{pump}   & $N_{\text{pump}}$  & pump axis $\omega_1$ \\
    \code{probe}  & $N_{\text{probe}}$ & probe axis $\omega_3$ \\
    \code{delays} & $N_{t_2}$ & population times $t_2$ \\
    \code{units}  & dict & unit labels \\
    \bottomrule
  \end{tabular}
  \caption{Fields stored on \code{Dataset2D}.}
  \label{tab:fields2d}
\end{table}

\subsection{Loading and the loader registry}
As in 1-D, formats are resolved through a registry
(\code{pymorgan.twoD.load}). The registered readers are \code{P2DAT} (processed
maps), \code{MESS\_2DIR}, \code{UoS\_2DIR}, \code{RAL\_RAW} (raw
interferograms, phased and transformed on load) and \code{RAL\_Proc};
\code{pm.available\_map\_loaders()} lists what the running installation
recognises.

\begin{lstlisting}
import pymorgan as pm

data = pm.load_2D("scan.p2dat", data_type="P2DAT")
print(pm.available_map_loaders())
\end{lstlisting}

A new format is registered with \code{@pm.register\_map\_loader("Name")}; the
reader must return \code{(Z, pump, probe, delays, units, freq\_units)}.

\paragraph{The P2DAT format.}
The first row is \code{0, 0, $t_2^{(1)}$, $t_2^{(2)}$, \dots} (the population
times); each subsequent row is \code{pump, probe, S($t_2^{(1)}$), \dots}. The
flat pump/probe ordering is reshaped (Fortran order) into the
\code{Z[pump, probe, $t_2$]} cube. Data is in m\,OD.

\subsection{Background correction}
The 2-D background is a chosen $t_2$ map subtracted from the others,
\begin{equation}
  Z_C(\omega_1,\omega_3,t_2) = Z(\omega_1,\omega_3,t_2) - Z(\omega_1,\omega_3,t_2^{\text{ref}}),
\end{equation}
with the reference map itself left unchanged so it can still be inspected:

\begin{lstlisting}
data.background_correct(reference=-1)   # by t2 index
data.background_correct(t2=50)          # by nearest t2 value
data.background_correct(do_correct=False)  # passthrough
\end{lstlisting}

\subsection{Plotting}
\code{plot\_map} renders the map nearest a requested $t_2$ as a pump-vs-probe
contour and returns a \code{Map2DAxes(ax, top, cbar)} named tuple. It is built
\emph{from} the supplied (or freshly created) axis: the colorbar and an optional
top panel are appended with a divider, so the whole composite drops into a
parent layout.

\begin{lstlisting}
out = data.plot_map(0.5)                 # one t2 map; new figure by default
out = data.plot_map(0.5, ax=my_ax, show_colorbar=False)   # into a layout
\end{lstlisting}

Two behaviours are specific to 2-D:
\begin{itemize}
  \item \textbf{Diagonal.} The $\omega_1 = \omega_3$ line is drawn only where the
        pump and probe axes overlap (pump and probe can be set independently);
        when their ranges are disjoint no diagonal is shown.
  \item \textbf{Top spectrum.} A steady-state or pump--probe spectrum can be drawn
        in a panel above the map (sharing the pump axis) via \code{top\_spectrum},
        which accepts a \code{Spectrum}, an \code{\{"X","Y"\}} mapping, or an
        \code{(x, y)} pair:
\end{itemize}

\begin{lstlisting}
ftir = pm.load_spectrum("sample_FTIR.csv", kind="absorption")
data.plot_map(0.5, top_spectrum=ftir, top_label="FTIR")
\end{lstlisting}

The same \code{show\_xlabel} / \code{show\_ylabel} / \code{show\_colorbar} /
\code{show\_colorbar\_label} toggles as in 1-D apply. The frequency-axis
labelling convention is the \code{freq\_label} setting
(\S\ref{sec:settings}): \code{Omega\_n} ($\omega_1/\omega_3$), \code{Omega\_PP}
($\omega_{\text{pump}}/\omega_{\text{probe}}$), \code{Pump-Probe}, or
\code{Omega\_n/2pic}; it can be overridden per call.

\subsubsection{Axis units and the $10^3$ rule}\label{sec:twoD:units}
The maps are stored in wavenumbers, but the axes can be displayed in any of the
units of Table~\ref{tab:2dunits}, selected with the \code{twoD\_freq\_unit}
setting or the \code{freq\_unit} argument of \code{plot\_map} and
\code{plot\_surface}. Conversions pivot through the wavelength,
\begin{equation}
  \lambda\,[\mathrm{nm}] = \frac{P_{\text{in}}}{x_{\text{in}}},
  \qquad
  x_{\text{out}} = \frac{P_{\text{out}}}{\lambda\,[\mathrm{nm}]},
\end{equation}
with the constants $P$ listed in the table; wavenumbers in reciprocal inches
follow $\tilde\nu\,[\mathrm{in}^{-1}] = 2.54\,\tilde\nu\,[\mathrm{cm}^{-1}]$.

\begin{table}[h]
  \centering
  \begin{tabular}{llll}
    \toprule
    Token & Quantity & $P$ (unit\,$\cdot$\,nm) & Typical mid-IR value \\
    \midrule
    \code{cm-1} & wavenumber & $10^{7}$              & 2000 \\
    \code{in-1} & wavenumber & $2.54\times10^{7}$    & 5080 \\
    \code{nm}   & wavelength & ---                   & 5000 \\
    \code{THz}  & frequency  & $2.99792458\times10^{5}$ & 60 \\
    \code{eV}   & energy     & $1239.841984$         & 0.25 \\
    \bottomrule
  \end{tabular}
  \caption{Display units of the 2-D spectral axes.}
  \label{tab:2dunits}
\end{table}

An axis whose values exceed $5\times10^{3}$ in the display unit is additionally
divided by $1000$, and the label gains a $10^{3}$ prefix: a visible-range map
reads $\omega_1$ ($10^{3}$ cm$^{-1}$) with ticks at $18\ldots22$ rather than
five-digit numbers. Axis limits and the CLS/IvCLS/NLS overlays are converted
together with the axes, so a region selected in wavenumbers stays on the same
features after a unit change.

\subsection{Analysis}
The analyse stage (\code{pymorgan.twoD.analyse}) offers map and diagonal
extraction --- the pump\,=\,probe diagonal is defined over the overlap of the two
axes --- together with the spectral-diffusion metrics of
\S\ref{sec:twoD:sd} and the Kubo lineshape fitting of
\code{pymorgan.twoD.kubo\_fit}.

\begin{lstlisting}
pump, probe, m = data.slice_at(0.5)   # the 2-D map at this t2
freq, sig      = data.diagonal(0.5)   # signal along omega_1 = omega_3

cls  = data.center_line_slope()              # [t2, slope]
both = data.center_line_slope(both=True)     # [t2, CLS, IvCLS]
nls  = data.nodal_line_slope()               # [t2, slope]
\end{lstlisting}

\subsection{Spectral diffusion}\label{sec:twoD:sd}
The centre-line slope (CLS) is obtained by taking, for every pump frequency
$\omega_1$, the extremum of the slice along the probe axis and fitting a line
through those points; the inverse CLS (IvCLS) swaps the roles of the two axes,
and the nodal-line slope (NLS) tracks the zero crossing between the bleach and
the excited-state absorption. The line is fitted robustly (Huber loss), so a few
outlying slices do not tilt the result.

\subsubsection{Sub-pixel refinement}
The extremum of a slice is \emph{not} taken at the sampled maximum: with probe
pixels of a few cm$^{-1}$ that would quantise the centre line and bias the
slope. Instead the position is obtained in two stages:
\begin{enumerate}
  \item the slice is interpolated onto a grid \code{sd\_interpolation\_factor}
        times finer with a cubic spline, and the extremum is located there,
        already resolving positions below the spectral resolution;
  \item a local model --- \code{sd\_interpolation}: a quadratic, cubic or
        quartic polynomial, or a Gaussian or Lorentzian peak --- is fitted by
        least squares to the \emph{measured} points around that location, and
        its stationary point is the reported centre.
\end{enumerate}
Fitting the measured samples (rather than applying a three-point formula to the
interpolated curve) keeps the estimate continuous and averages the noise over
several points. On a synthetic tilted trough sampled at 2.5\,cm$^{-1}$ the
residual error is $\approx0.01$ pixel for the quadratic model and
$\approx0.003$ pixel for the quartic one; with \code{sd\_interpolation = "none"}
the fine-grid position is returned unrefined, and its resolution is set by the
interpolation factor alone.

Two further settings restrict what enters the fit:
\code{sd\_peak\_range} (cm$^{-1}$) is the half-width of the search window around
the coarse extremum, \code{sd\_intensity\_threshold} the minimum amplitude
(relative to the strongest signal of the map) a slice must reach to contribute,
and \code{sd\_peak\_type} selects the bleach minimum or the excited-state
absorption maximum.

The decay of the extracted slopes against $t_2$ is plotted from the GUI with
single- or bi-exponential fits. That figure always uses a \emph{linear} $t_2$
axis: population times are a short, linearly sampled series, unlike the wide
delay range of a 1-D experiment.

\subsection{Worked example}
A complete script is provided in \code{examples/two\_d\_2DIR.py}:

\begin{lstlisting}
import pymorgan as pm

pm.load_settings("settings.toml")
pm.apply_style()

data = pm.load_2D("testData/test2DIR.p2dat", data_type="P2DAT")  # load
data.background_correct(reference=-1)                             # process
data.plot_map(0.5, ShowLines=True)                               # plot
\end{lstlisting}

\subsection{Graphical interface}\label{sec:twoD:gui}

The GUI (\code{pymorgan-gui}) exposes the full 2-D pipeline through the
\textbf{2D} tab.  The left column contains a hierarchical file browser,
format selector, and a series of sub-tabs that drive processing and
analysis.  Directories shown in grey (60\,\% grey) are navigation
placeholders (\code{[..]} / \code{[...]}) that contain no loadable data.

\subsubsection{Phasing/FT sub-tab}

Raw (interferometer) 2-D data are Fourier-transformed and phased through
the \emph{Phasing/FT} sub-tab before any spectral map is rendered.
Table~\ref{tab:phasing} summarises the available controls.

\begin{table}[h]
  \centering
  \begin{tabular}{lll}
    \toprule
    Control & Options & Default \\
    \midrule
    Apodisation & Box, \textbf{Cos}, Cos$^2$, Cos$^3$, Hanning, Hamming, None & Cos \\
    Phase fit   & Constant, \textbf{Linear}, Quadratic, Cubic & Linear \\
    Phase range & $\pm N$ frequency units around peak centre & 10 \\
    Zero-pad    & enabled / disabled, factor, next-2$^k$ & disabled \\
    Pump correction & on / off & off \\
    \bottomrule
  \end{tabular}
  \caption{Phasing/FT controls and their defaults.}
  \label{tab:phasing}
\end{table}

The \textbf{Phase coeffs.}\ field (read-only) displays the polynomial
coefficients of the fitted phase for the currently selected delay,
formatted as \texttt{\%.3g}.  It updates automatically when the delay
changes or after a new load.

\subsubsection{Other plt. sub-tab}

The \emph{Other plt.} sub-tab contains three diagnostic views selected
by the \textsc{td} / \textsc{ph} / \textsc{cal} toggle buttons:

\begin{description}
  \item[TD] Time-domain interferogram and signal for the selected
    probe pixel.
  \item[PH] Fourier transform amplitude (black), a Gaussian fit
    (green dashed), and — for interferometer data — the unwrapped
    phase (red dots) and the polynomial fit (blue line) on a twin
    $y$-axis.  The $x$-axis is locked to $\pm 2\times\mathrm{FWHM}$
    of the Gaussian fit.  The phase axis is re-centred so the fitted
    phase is zero at the spectral peak, wrapped to $[-\pi,\pi]$, and
    plotted with fixed limits $[-\pi,\pi]$.
  \item[CAL] Calibrated probe axis (pixel vs.\ wavenumber / THz).
\end{description}

Clicking the \textbf{Pop} button opens an independent, resizable window
that mirrors the same axes.  The window is linked to the main GUI: it
updates whenever the delay or plot mode changes.  Closing the pop-out
does not affect the embedded plot.

\subsubsection{Movie / GIF export dialog}

The \textbf{Make Movie\ldots} button opens \code{movie\_dialog.py}, which
exports an animated GIF or MP4 video of the 2-D contour map stepping
through all population times~$t_2$. Controls include:

\begin{description}
  \item[Delay range] Lower and upper $t_2$ limits to include in the animation.
  \item[Frame rate / duration] Frames per second or per-frame hold time in ms.
  \item[Figure size] Figure width and height in inches.
  \item[Colour scale ($Z$-scale)] Fixed or per-frame normalised colour range.
  \item[Output format] GIF (via Pillow) or MP4 (via the matplotlib
        \code{FFMpegWriter}; requires \code{ffmpeg} on \code{PATH}).
\end{description}

A progress bar reports frame rendering. The dialog uses the plot-controls
panel settings (CLS/IvCLS overlays, quick-plot mode, etc.) that are active
in the main window when the dialog is opened.

% PyMORGAN — Steady-state spectra chapter
\section{Steady-state spectra}\label{sec:steadyState}

The \code{pymorgan.steadyState} module provides lightweight wrappers for
absorption, emission and excitation spectra. Its design mirrors the
one-dimensional branch (chapter~\ref{sec:oneD}): data are held in a plain
object (\code{Spectrum}), file formats are resolved through a loader registry,
and presentation is driven by the shared \code{Settings}.

\subsection{Data model}

\subsubsection{\code{SpectrumKind}}

\code{SpectrumKind} is a string enum that labels the physical quantity
measured:

\begin{lstlisting}
class SpectrumKind(StrEnum):
    ABSORPTION = "absorption"
    EMISSION   = "emission"
    EXCITATION = "excitation"
\end{lstlisting}

\subsubsection{\code{Spectrum}}

A single steady-state measurement. The spectral axis \code{x} and signal
\code{y} are stored as \code{numpy} arrays; the remaining fields are metadata:

\begin{table}[h]
  \centering
  \begin{tabular}{lll}
    \toprule
    Attribute & Type & Meaning \\
    \midrule
    \code{x}       & \code{ndarray} & Spectral axis (values in \code{x\_units}) \\
    \code{y}       & \code{ndarray} & Signal (absorbance or intensity) \\
    \code{kind}    & \code{SpectrumKind} & Measurement type \\
    \code{x\_units}& \code{str} & Spectral-axis unit token (\code{"nm"}, \code{"cm-1"}, \code{"eV"}) \\
    \code{label}   & \code{str | None} & Legend label \\
    \code{source}  & \code{str | None} & File path (set automatically on load) \\
    \bottomrule
  \end{tabular}
  \caption{Attributes of a \code{Spectrum} object.}
  \label{tab:spectrum:fields}
\end{table}

\subsubsection{\code{SpectrumSeries}}

An ordered list of \code{Spectrum} objects for overlay plots. It supports
iteration, indexing and the standard sequence protocol. Construction from
multiple files returns a \code{SpectrumSeries} directly.

\subsection{Loading and the loader registry}

The module exposes a loader registry analogous to the 1-D and 2-D registries.
The default \code{"csv"} reader handles comma- and tab-delimited files with
automatic delimiter detection. Custom readers are registered with the
\code{@pm.register\_spectrum\_loader} decorator:\

\begin{lstlisting}
import pymorgan as pm

# List available loaders
print(pm.available_spectrum_loaders())

# Load a single spectrum
sp = pm.load_spectrum("sample.csv", kind="absorption")

# Specify units explicitly (overrides the reader's guess)
sp = pm.load_spectrum("sample.csv", kind="emission", x_units="nm", label="Sample A")

# Load directly via Dataset1D convenience method
sp2 = pm.Spectrum.from_file("sample.opus", kind="absorption", data_type="opus")
\end{lstlisting}

The currently registered loaders are \code{csv} (generic delimiter-detected
text), \code{opus} (Bruker OPUS binary, via \code{load.py}), \code{uv\_vis}
(Perkin-Elmer / Shimadzu UV-Vis text export) and \code{fluorimeter} (Horiba
FluorEssence \code{.fes} files).

A custom loader returns a three-element tuple:
\begin{lstlisting}
@pm.register_spectrum_loader("MySpectrometer")
def read_my_spectrometer(path):
    ...
    return x, y, x_units    # x_units may be None (falls back to "nm")
\end{lstlisting}

\subsection{Transformations}

Two per-\code{Spectrum} transformations return a new \code{Spectrum}
(the original is unchanged):

\begin{description}
  \item[\code{normalised()}] Scales the signal so that
    $\max(|y|) = 1$.
  \item[\code{to\_stimulated\_emission()}] Computes the stimulated-emission
    cross-section $F_\text{stim} \propto \bar{\nu}^4 \, F_\text{spont}$ (valid
    for emission spectra on a wavelength axis). Useful when overlaying on
    transient-absorption spectra.
\end{description}

\code{SpectrumSeries.normalised()} applies the same operation to every member
and returns a new series.

\subsection{Integration with the 1-D pipeline}

A \code{Spectrum} can be passed directly to \code{Dataset1D.plot\_spectra} as a
steady-state overlay:\

\begin{lstlisting}
ftir = pm.load_spectrum("background_FTIR.csv", kind="absorption")
data.plot_spectra([0.5, 1, 5], ss_overlay=ftir)
\end{lstlisting}

\code{as\_overlay\_dict()} converts a \code{Spectrum} to the
\code{\{"X": x, "Y": y\}} mapping that \code{plot\_spectra} also accepts,
for cases where you want to pass raw arrays.

\subsection{Plotting}

Both \code{Spectrum} and \code{SpectrumSeries} have a \code{plot()} method.
The axis-label convention and delimiter follow the active settings
(\S\ref{sec:settings}), and optional keyword arguments are forwarded to
\code{matplotlib.Axes.plot}.

\begin{lstlisting}
import pymorgan as pm

pm.load_settings("settings.toml")
pm.apply_style()

# --- Single spectrum ---
sp = pm.load_spectrum("sample.csv", kind="absorption", label="Sample")
sp.plot()

# --- Normalised emission ---
em = pm.load_spectrum("emission.csv", kind="emission", label="Emission")
em.normalised().plot()

# --- Series: multiple spectra coloured along a rainbow colourmap ---
series = pm.SpectrumSeries.from_files(
    ["a.csv", "b.csv", "c.csv"],
    kind="absorption",
    labels=["10 uM", "50 uM", "100 uM"],
)
series.plot(cmap="viridis", normalise=True)

pm.show_plots()
\end{lstlisting}

The per-call keyword arguments are identical to those of the 1-D plotters:
\code{ax=} to draw into an existing axis, \code{label\_style=} to override the
delimiter, and \code{normalise=} to scale the signal to unit peak.


\section{Settings}\label{sec:settings}

Every presentation choice that is not intrinsic to the data lives in a single
\code{Settings} object: the matplotlib style profile, label and unit
conventions, colour scale, axis units and the parameters of the spectral
diffusion analysis. One instance is active at a time; the plotting routines read
it, and any call may override individual values.

\begin{lstlisting}
import pymorgan as pm

pm.load_settings("settings.toml")     # read a file
pm.update_settings(cmap="Seismic", quick_plots=True)
s = pm.get_settings()                 # the active object
pm.save_settings("settings.toml")     # write it back, comments preserved

with pm.use_settings(profile="poster"):   # temporary override
    data.plot_contour()
\end{lstlisting}

Values are coerced to the declared type, so a string from a text field or a TOML
file becomes the right enum, number, path or pair: \code{pm.update\_settings(
contour\_figsize="9, 5")} and \code{pm.update\_settings(x\_axis\_unit="nm")} both
work. An unknown name raises \code{TypeError} rather than being silently stored.

\subsection{The configuration file}
\code{settings.toml} is read at start-up (the GUI looks in the working directory
first, then next to the package) and written back by \emph{Save permanently} in
the settings dialog. The file is edited in place, so comments and ordering
survive a round trip. Keys map one-to-one onto the tables below; a key that is
absent keeps its default, and an unknown key is ignored.

\subsection{Editing settings in the GUI}
\emph{View $\to$ Aesthetics / Settings\ldots} opens a panel whose widgets are
generated from the settings themselves and grouped in four tabs --- Common, 1D,
2D, and GUI \& Defaults --- matching the tables that follow. Changes apply to
the current session immediately (the embedded preview re-renders);
\emph{Save permanently} writes them to \code{settings.toml}.

\subsection{Common settings}
Shared by every plot: style profile, labels, colour scale and the units of the spectral axis.

\begin{table}[h]
  \centering
  \small
  \begin{tabular}{p{0.26\textwidth}p{0.28\textwidth}p{0.30\textwidth}l}
    \toprule
    Key & Label in the panel & Type & Default \\
    \midrule
    \code{profile} & Style profile & choice: \code{regular}, \code{poster}, \code{inset} & \code{regular} \\
    \code{font\_scale} & Font scale & float & \code{1} \\
    \code{label\_style} & Axis-label style & choice: \code{()}, \code{[]}, \code{/}, (empty) & \code{()} \\
    \code{delta\_a\_units} & \dA{} units & choice: \code{mOD}, \code{x1E3} & \code{x1E3} \\
    \code{cmap} & Colourmap & choice: \code{DkRd/Wh/DkBu}, \code{Rd/Wh/Bu v2}, \code{Seismic}, \code{Jet}, \code{vik}, \code{berlin} & \code{DkRd/Wh/DkBu} \\
    \code{white\_levels} & White/Zero levels & float & \code{0} \\
    \code{asinh\_pct} & arcsinh linear \% & float & \code{5.0} \\
    \code{x\_axis\_unit} & Spectral X-axis unit & choice: \code{native}, \code{nm}, \code{cm-1}, \code{eV}, \code{THz} & \code{native} \\
    \code{secondary\_axis} & Complementary X axis & bool & \code{false} \\
    \code{round\_labels} & Round trace labels & bool & \code{true} \\
    \code{quick\_plots} & Use quick plots (fast redraws) & bool & \code{false} \\
    \code{ss\_line\_dashes} & Steady-state dash pattern (on, off) & text & \code{4, 3} \\
    \bottomrule
  \end{tabular}
  \caption{Common settings.}
  \label{tab:settings:common}
\end{table}

\subsection{1-D settings}
Delay-axis scaling and the appearance of kinetic / transient-spectra traces.

\begin{table}[h]
  \centering
  \small
  \begin{tabular}{p{0.26\textwidth}p{0.28\textwidth}p{0.30\textwidth}l}
    \toprule
    Key & Label in the panel & Type & Default \\
    \midrule
    \code{time\_axis\_scale} & Time-axis scale & choice: \code{symlog}, \code{log}, \code{lin} & \code{symlog} \\
    \code{time\_axis\_label} & Time-axis labels & choice: \code{power}, \code{decimal}, \code{mixed} & \code{power} \\
    \code{kinetics\_marker\_size} & Kinetics marker size & float & \code{6} \\
    \code{kinetics\_line\_width} & Kinetics line width & float & \code{1.5} \\
    \code{legend\_label\_spacing} & Legend entry spacing & float & \code{0.5} \\
    \code{prune\_legend} & Prune legend & bool & \code{true} \\
    \code{max\_legend\_entries} & Max legend entries & float & \code{15} \\
    \code{ss\_abs\_color} & Steady-state abs. colour & text & \code{b} \\
    \code{ss\_em\_color} & Steady-state em. colour & text & \code{r} \\
    \code{ss\_line\_width} & Steady-state line width & float & \code{1.5} \\
    \code{ss\_line\_alpha} & Steady-state line alpha & float & \code{0.15} \\
    \code{ss\_fill\_alpha} & Steady-state fill alpha & float & \code{0.05} \\
    \code{inherit\_cut\_limits} & Cuts inherit plot limits & bool & \code{true} \\
    \code{load\_single\_scans} & Load single scans & bool & \code{false} \\
    \code{chirp\_boundary\_handling} & Chirp boundary handling & choice: \code{crop}, \code{extrapolate} & \code{crop} \\
    \code{parallel\_fitting} & Chirp parallel fitting & choice: \code{threadpool}, \code{processpool}, \code{disabled} & \code{threadpool} \\
    \code{solvent\_per\_pixel\_dt} & Per-pixel solvent t$_0$ offset & bool & \code{false} \\
    \code{solvent\_per\_pixel\_scale} & Per-pixel solvent amplitude scale & bool & \code{true} \\
    \code{error\_shading} & Error shading ($\pm$1 std) & bool & \code{false} \\
    \code{error\_shading\_alpha} & Error shading alpha & float & \code{0.2} \\
    \bottomrule
  \end{tabular}
  \caption{1-D settings.}
  \label{tab:settings:oneD}
\end{table}

\subsection{2-D settings}
Map appearance, axis units and the spectral-diffusion analysis.

\begin{table}[h]
  \centering
  \small
  \begin{tabular}{p{0.26\textwidth}p{0.28\textwidth}p{0.30\textwidth}l}
    \toprule
    Key & Label in the panel & Type & Default \\
    \midrule
    \code{freq\_label} & 2-D frequency labels & choice: \code{Omega\_n}, \code{Omega\_PP}, \code{Pump-Probe}, \code{Omega\_n/2pic} & \code{Omega\_n} \\
    \code{pump\_axis} & 2-D pump axis orientation & choice: \code{Horizontal}, \code{Vertical} & \code{Horizontal} \\
    \code{twoD\_freq\_unit} & 2-D axis unit & choice: \code{cm-1}, \code{in-1}, \code{nm}, \code{THz}, \code{eV} & \code{cm-1} \\
    \code{t2\_label\_style} & t$_2$ delay label format & choice: \code{t2=}, \code{plain} & \code{t2=} \\
    \code{cls\_color} & CLS colour & text & \code{b} \\
    \code{ivcls\_color} & IvCLS colour & text & \code{r} \\
    \code{nls\_color} & NLS colour & text & \code{g} \\
    \code{overlay\_draw\_style} & Overlay draw style & choice: \code{Points}, \code{Lines}, \code{Both} & \code{Points} \\
    \code{overlay\_linewidth} & Overlay linewidth & float & \code{1.5} \\
    \code{sd\_intensity\_threshold} & Intensity threshold (top x\%) & float & \code{0.1} \\
    \code{sd\_peak\_range} & Peak search range (cm$^{-1}$) & float & \code{10} \\
    \code{sd\_peak\_type} & Target peak type & choice: \code{min}, \code{max} & \code{min} \\
    \code{sd\_interpolation} & Peak interpolation & choice: \code{quadratic}, \code{cubic}, \code{quartic}, \code{gaussian}, \code{lorentzian}, \code{none} & \code{quadratic} \\
    \code{sd\_interpolation\_factor} & Interpolation factor (axis upsampling) & int & \code{4} \\
    \code{kubo\_ftir\_weight} & Kubo joint fit FTIR vs 2D weight & float & \code{1} \\
    \bottomrule
  \end{tabular}
  \caption{2-D settings.}
  \label{tab:settings:twoD}
\end{table}

\subsection{GUI and defaults settings}
Application preferences and the values pre-selected when the GUI starts.

\begin{table}[h]
  \centering
  \small
  \begin{tabular}{p{0.26\textwidth}p{0.28\textwidth}p{0.30\textwidth}l}
    \toprule
    Key & Label in the panel & Type & Default \\
    \midrule
    \code{gui\_theme} & GUI theme & choice: \code{Fusion}, \code{Fusion Dark} & \code{Fusion} \\
    \code{splash\_duration} & Splash-screen duration (s, 0 disables) & float & \code{2} \\
    \code{default\_datadir} & Default data directory & text & (none) \\
    \code{default\_oneD\_datatype} & Default 1D data type & choice: \code{Helios\_TA}, \code{MESS\_TRIR}, \code{MESS\_TRUVIS}, \code{PDAT}, \code{UniGE\_FLUPSnew}, \code{UniGE\_FLUPSold}, \code{UniGE\_fsTA}, \code{UniGE\_nsTA}, \code{UoS\_IRpp} & \code{PDAT} \\
    \code{default\_twoD\_datatype} & Default 2D data type & choice: \code{MESS\_2DIR}, \code{P2DAT}, \code{RAL\_Proc}, \code{RAL\_RAW}, \code{UoS\_2DIR} & \code{P2DAT} \\
    \code{default\_calibration\_datatype} & Default calibration setup & text & \code{UniGE TRIR (Intensity)} \\
    \code{styles\_dir} & Style directory (blank = bundled) & text & (none) \\
    \code{contour\_figsize} & Contour figure size (w, h in) & text & \code{8, 6} \\
    \code{kinetics\_figsize} & Kinetics figure size (w, h in) & text & \code{7.5, 4.5} \\
    \code{spectra\_figsize} & Spectra figure size (w, h in) & text & \code{7.5, 4.375} \\
    \bottomrule
  \end{tabular}
  \caption{GUI and defaults settings.}
  \label{tab:settings:gui}
\end{table}

\subsection{Declared axis units}\label{sec:units:axes}
Both dataset classes report the units of their axes through
\code{axis\_units()}, which returns a \code{DatasetUnits} record holding one
entry per axis: \code{x} and \code{y} for the plotted axes, \code{z} for the
signal, and \code{t2} for the population time of a 2-D dataset (\code{None} in
1-D, where the delay is already a plotted axis). Each entry carries the measured
quantity, the machine-readable unit token, the mathtext used on plots, and the
factor the values are divided by for display.

\begin{lstlisting}
u = data.axis_units()            # 1-D: x = probe, y = delay, z = signal
u.x.unit                         # 'cm-1'
u.x.label()                      # 'Wavenumber (cm$^{-1}$)'
u.y.quantity                     # 'Delay'
print(u)                         # a readable summary of all axes
u.as_dict()                      # plain nested dict, e.g. for export

v = d2.axis_units()              # 2-D: x = pump, y = probe, z = signal, t2 = delay
v.t2.unit                        # 'ps'
\end{lstlisting}

By default the \emph{stored} units are reported. With \code{display=True} the
axes are described as they would be plotted: converted to
\code{x\_axis\_unit} (1-D) or \code{twoD\_freq\_unit} (2-D), carrying the
$10^3$ factor where one is applied (\S\ref{sec:twoD:units}), and with \code{x}
and \code{y} swapped when \code{pump\_axis} is \code{Vertical}. The
\code{label()} method renders the axis label in any of the delimiter
conventions of \code{label\_style}, so a caller can annotate a figure of its own
without duplicating the unit logic.

\subsection{Notes on individual settings}
\begin{description}
  \item[\code{quick\_plots}] Draws contour maps as a single image
    (\code{pcolormesh}) and skips the contour lines, which makes dragging the
    colour-scale slider or stepping through $t_2$ far quicker. It overrides the
    \emph{Filled} and \emph{Show lines} options of the plot-controls panel, so
    turn it off for final figures.
  \item[\code{x\_axis\_unit}, \code{secondary\_axis}] Control the spectral axis
    of 1-D plots: the axis is converted to the chosen unit, and a complementary
    axis (cm$^{-1}\leftrightarrow$nm) can be added on the opposite side.
  \item[\code{twoD\_freq\_unit}] The same idea for the pump/probe axes of 2-D
    maps, including reciprocal inches, with the $10^3$ rescaling described in
    \S\ref{sec:twoD:units}.
  \item[\code{sd\_*}] The spectral-diffusion family: search window, amplitude
    threshold, peak type, and the interpolation factor and model used for the
    sub-pixel refinement (\S\ref{sec:twoD:sd}).
  \item[\code{parallel\_fitting}] Chooses how the per-pixel chirp fit is
    distributed: a thread pool, a process pool (true parallelism, best for large
    datasets), or sequentially on the calling thread.
  \item[\code{gui\_theme}] Any Qt style available on the machine; appending
    \emph{Dark} (e.g. \code{Fusion Dark}) switches to the dark palette, which
    also recolours the embedded plots.
\end{description}

% PyMORGAN — Calibration chapter
\section{Spectrometer \& Pulse Shaper Wavelength Calibration}\label{sec:calibration}

Wavelength and wavenumber axis calibration maps detector pixel indices $p \in \{1, \dots, N_\text{pix}\}$ to physical spectral values (wavelength $\lambda$ in nm or wavenumber $\tilde{\nu}$ in \text{cm}$^{-1}$). PyMORGAN provides a dedicated calibration engine in the \code{pymorgan.cal} module, supporting both \emph{spectrometer calibration} (for probe beam spectrographs) and \emph{pulse shaper calibration} (for optical pulse shaper pump beam masks).

\subsection{Pipeline Overview}

The calibration workflow matches an experimentally measured absorption spectrum $I_\text{meas}(p)$ against a high-resolution reference FTIR spectrum $A_\text{ref}(\lambda)$ (such as liquid 1,4-dioxane, liquid polystyrene, or atmospheric water vapour). Non-linear least-squares (NLS) optimisation fits optical dispersion parameters so that the synthetic spectrum evaluated through the dispersion function reproduces the reference absorption features.

The workflow proceeds in six stages:
\begin{enumerate}
  \item \textbf{Load Reference Spectrum:} Load $A_\text{ref}(\lambda)$ or $A_\text{ref}(\tilde{\nu})$ with automatic delimiter and unit detection (\code{load\_reference\_spectrum}).
  \item \textbf{Load Experimental Data:} Parse raw detector intensity or absorbance arrays for one of ten hardware setups (\code{load\_experimental\_spectrum}).
  \item \textbf{Pre-alignment (Template Scanning):} Perform 1D cross-correlation scanning to estimate initial wavelength shift $\lambda_1$ at pixel~1.
  \item \textbf{NLS Parameter Fitting:} Optimise dispersion coefficients and additive polynomial baseline using \code{scipy.optimise.least\_squares} with 1st or 2nd derivative filtering.
  \item \textbf{Axis Inversion:} Invert pixel-to-wavelength mapping $p(\lambda) \to \lambda(p)$ analytically or numerically.
  \item \textbf{Export \& Integration:} Save calibrated axis arrays (\code{CalibratedProbe.csv}, \code{CalibratedPump.csv}, \code{SingleMask.txt}, \code{MultipleMask.txt}).
\end{enumerate}

\subsection{Data Model \& Data Structures}

The \code{pymorgan.cal} package defines four primary immutable data structures:

\begin{description}
  \item[\code{ReferenceSpectrum}:] Container for reference spectra library items.
\begin{lstlisting}
class ReferenceSpectrum(NamedTuple):
    name: str             # Spectrum name (e.g. 'FTIR-Dioxane')
    spectral_axis: ndarray# Spectral grid (nm or cm-1)
    absorbance: ndarray   # Absorbance array A(x)
    axis_unit: str        # Unit hint ('nm' or 'cm-1')
\end{lstlisting}

  \item[\code{ExperimentalData}:] Container for loaded experimental spectrograph measurements.
\begin{lstlisting}
class ExperimentalData(NamedTuple):
    detector_data: list[ndarray]  # Raw detector intensity/absorbance arrays
    gratings: list[float]         # Grating groove density (grooves/mm)
    cwl: list[float]              # Central wavelength setting (nm or cm-1)
    file_path: Path               # File path on disk
    min_wl: float | None = None   # Lower spectral bound from metadata
    max_wl: float | None = None   # Upper spectral bound from metadata
\end{lstlisting}

  \item[\code{CalibrationResult}:] Output container returned by \code{fit\_wavelength\_axis}.
\begin{lstlisting}
class CalibrationResult(NamedTuple):
    wavelength_nm: ndarray   # Calibrated wavelength axis (nm)
    wavenumber_cm1: ndarray  # Calibrated wavenumber axis (cm-1)
    fit_params: ndarray      # Optimised dispersion parameters vector p
    residuals: ndarray       # Fit residual array
    ref_x_cut: ndarray       # Reference spectral axis in active cut region
    ref_y_cut: ndarray       # Reference absorbance in active cut region
    model_y_fit: ndarray     # Modelled experimental absorbance fit
    baseline_fit: ndarray    # Additive polynomial baseline curve
    exp_corr_y_fit: ndarray  # Baseline-corrected experimental absorbance
    ref_y_corr: ndarray      # Baseline-corrected reference spectrum
\end{lstlisting}

  \item[\code{SpectrumFitResult}:] Container for 1D Gaussian pump beam profile diagnostics.
\begin{lstlisting}
class SpectrumFitResult(NamedTuple):
    amplitude: float   # Peak amplitude (Counts)
    w0: float          # Central frequency w0 (cm-1 or nm)
    fwhm: float        # Full-Width at Half-Maximum (FWHM)
    fit_curve: ndarray # High-density fitted Gaussian profile
    axis_fit: ndarray  # High-density evaluation axis
\end{lstlisting}

  \item[\code{ProbeFitResult}:] Container returned by \code{fit\_probe\_spectrum}, which
  wraps up multi-detector Gaussian diagnostics for the probe beam profile.
\begin{lstlisting}
class ProbeFitResult(NamedTuple):
    fits: list[SpectrumFitResult]  # Per-detector Gaussian fit results
    detector_data: list[ndarray]   # Raw detector arrays that were fitted
\end{lstlisting}
\end{description}

\subsection{Mathematical Optical Dispersion Models}

PyMORGAN supports three physical models for mapping reference wavelength $\lambda$ (in nm) to detector pixel index $p(\lambda) \in [1, N_\text{pix}]$:

\subsubsection{Linear Grating Dispersion Model}
For diffraction grating spectrographs in first-order paraxial approximation, the pixel position $p(\lambda)$ varies linearly with wavelength around the central wavelength $\lambda_0$:
\begin{equation}\label{eq:dispersion_linear}
  p(\lambda) = (\lambda - \lambda_0) \cdot d + c_\text{pix}
\end{equation}
where $\lambda_0$ is the central wavelength mapped to the central pixel $c_\text{pix} = (1 + N_\text{pix})/2$, and $d$ is the linear dispersion in pixels/nm.

\subsubsection{Polynomial Grating Dispersion Model}
For wide-bandwidth spectrographs or non-paraxial grating setups (such as UniGE TRUVIS-II), higher-order geometric distortion is modelled via a 2nd- or 3rd-order polynomial expansion around central wavelength $\lambda_0$:
\begin{equation}\label{eq:dispersion_poly}
  p(\lambda) = (\lambda - \lambda_0) \cdot d + 10^{-4} \, c_2 \, (\lambda - \lambda_0)^2 + 10^{-7} \, c_3 \, (\lambda - \lambda_0)^3 + c_\text{pix}
\end{equation}
where $c_2$ and $c_3$ are quadratic and cubic dispersion curvature coefficients.

\subsubsection{Prism Dispersion Model (SF10 Glass)}
For prism-based spectrographs (e.g.\ UniGE NIR-TA using an SF10 Flint glass equilateral prism), dispersion is governed by Snell's law and the Sellmeier equation. The refractive index $n(\lambda)$ for SF10 glass is:
\begin{equation}\label{eq:sellmeier}
  n(\lambda) = \sqrt{1 + \sum_{i=1}^{3} \frac{B_i \, \lambda_{\mu\text{m}}^2}{\lambda_{\mu\text{m}}^2 - C_i}}
\end{equation}
with Sellmeier coefficients:
\begin{equation*}
  \begin{aligned}
    B_1 &= 1.62153902, & B_2 &= 0.256287842, & B_3 &= 1.64447552 \\
    C_1 &= 0.0122241457, & C_2 &= 0.0595736775, & C_3 &= 147.468793
  \end{aligned}
\end{equation*}
where $\lambda_{\mu\text{m}} = \lambda / 1000$. The deflection output angle $\theta_\text{out}(\lambda)$ from a prism with apex angle $\alpha = 60.84^\circ$ at incidence angle $\theta_\text{in}$ is:
\begin{equation}\label{eq:prism_angle}
  \theta_\text{out}(\lambda) = \arcsin\!\left[ n(\lambda) \sin\!\left( \alpha - \arcsin\!\left( \frac{\sin\theta_\text{in}}{n(\lambda)} \right) \right) \right]
\end{equation}
The pixel index $p(\lambda)$ is then determined by focal length $f_\text{mm}$ and detector pixel shift $x_\text{shift}$:
\begin{equation}\label{eq:prism_pixel}
  p(\lambda) = f_\text{mm} \cdot \left( \frac{1000}{25} \right) \cdot \sin\!\big[ \theta_\text{out}(\lambda) - \theta_\text{out}(\lambda_\text{ref}) \big] + x_\text{shift}
\end{equation}

\subsection{Forward Model, Baseline Correction, and Optimisation Cost Function}

\subsubsection{Forward Synthetic Spectrum Model}
Given measured experimental spectrum array $I_\text{meas}(p)$, the model synthesises an experimental absorbance curve $y_\text{model}(\lambda)$ matched to reference grid points $\lambda$:
\begin{equation}\label{eq:forward_model}
  y_\text{model}(\lambda) = I_\text{meas}\!\big[p(\lambda)\big] \cdot s + b_0 + 10^{-6} \, b_1 \, (\lambda - \lambda_0) + 10^{-6} \, b_2 \, (\lambda - \lambda_0)^2
\end{equation}
where $s$ is a global scaling factor, and $(b_0, b_1, b_2)$ parameterise a smooth background baseline polynomial.

\subsubsection{Derivative Cost Function}
To eliminate broad solvent/background baseline drift during fitting, optimisation is performed in derivative space using Savitzky--Golay numerical filters (window size $W$, polynomial order 2):
\begin{equation}\label{eq:residuals}
  \mathbf{r}(\mathbf{p}) = \frac{d^k}{d\lambda^k} y_\text{model}(\lambda; \mathbf{p}) - \frac{d^k}{d\lambda^k} A_\text{ref}(\lambda) \quad (k \in \{1, 2\})
\end{equation}
The parameters $\mathbf{p}$ are optimised by bounded Non-Linear Least Squares:
\begin{equation}\label{eq:least_squares}
  \min_{\mathbf{p}} \frac{1}{2} \|\mathbf{r}(\mathbf{p})\|_2^2 \quad \text{subject to } \mathbf{lb} \le \mathbf{p} \le \mathbf{ub}
\end{equation}

\subsection{Axis Inversion Procedures}

Once optimal parameters $\mathbf{p}^*$ are found, the calibrated wavelength axis $\lambda(p)$ for pixel indices $p \in \{1, \dots, N_\text{pix}\}$ is calculated:

\begin{itemize}
  \item \textbf{Linear Grating:} $\lambda(p) = \frac{p - 1}{d} + \lambda_1$.
  \item \textbf{Polynomial Grating:} Inverted via dense 1D grid interpolation over $\lambda \in [250, 1000]$\,nm.
  \item \textbf{Prism Model:} Inverted via KDTree spatial nearest-neighbor search against dense evaluation grid $p(\lambda)$.
\end{itemize}

Wavenumbers $\tilde{\nu}$ (in \text{cm}$^{-1}$) are directly converted via:
\begin{equation}
  \tilde{\nu}(p) = \frac{10^7}{\lambda(p)}
\end{equation}

\subsection{1D Gaussian Pump Spectrum Profile Diagnostics}

Pump beam profiles in optical pulse shaper calibration are fitted with a 3-parameter Gaussian distribution:
\begin{equation}\label{eq:gaussian_fit}
  G(x) = A \cdot \exp\!\left[ -4 \ln 2 \left( \frac{x - \omega_0}{\text{FWHM}} \right)^{\!2} \right]
\end{equation}
where $A$ is peak amplitude, $\omega_0$ is central frequency (\text{cm}$^{-1}$), and $\text{FWHM}$ is Full-Width at Half-Maximum.

\subsection{Supported Hardware Setup Configurations}

The calibration engine pre-configures ten spectrograph hardware setups:

\begin{table}[h]
  \centering
  \caption{Pre-configured calibration setup codes and default parameter estimates.}\label{tab:cal_setups}
  \begin{tabular}{llcccc}
    \toprule
    Code & Setup Name & Default $\lambda_0$ & Dispersion Guess & Degree & Default Bounds ($\Delta\lambda_\text{rel}$) \\
    \midrule
    1  & UoS TRIR                 & 2000 cm$^{-1}$ & 0.30 px/cm$^{-1}$ & 2      & [-1000, +1666] nm \\
    2  & UniGE TA                 & 530 nm         & 1.10 px/nm        & 2      & [-180, +210] nm \\
    3  & UniGE nsTA               & 530 nm         & 1.10 px/nm        & 2      & [-180, +210] nm \\
    4  & UZH Lab 2                & 2000 cm$^{-1}$ & 0.30 px/cm$^{-1}$ & 2      & [-455, +1250] nm \\
    5  & UniGE NIR-TA             & 1000 nm        & 1.00 px/nm        & Prism  & [-200, +600] nm \\
    6  & RAL Absorbance           & 2000 cm$^{-1}$ & 0.30 px/cm$^{-1}$ & 2      & [-653, +882] nm \\
    7  & RAL Intensity            & 2000 cm$^{-1}$ & 0.30 px/cm$^{-1}$ & 2      & [-653, +882] nm \\
    8  & UniGE TRIR (Int.)        & 2000 cm$^{-1}$ & 0.30 px/cm$^{-1}$ & 2      & [-200, +200] nm \\
    9  & UniGE TRIR (Abs.)        & 2000 cm$^{-1}$ & 0.30 px/cm$^{-1}$ & 2      & [-200, +200] nm \\
    10 & UniGE TRUVIS-II          & 530 nm         & 1.20 px/nm        & 3      & [-150, +220] nm \\
    \bottomrule
  \end{tabular}
\end{table}

The initial parameter bounds and dispersion guesses for each mode can be
queried programmatically with \code{get\_default\_calibration\_params(cal\_type\_code)},
which returns a dictionary with keys \code{p0}, \code{lb}, \code{ub},
\code{cwl\_guess} and \code{degree}.

\subsection{Graphical User Interface (Calibration Tab)}

The GUI incorporates a dedicated \textbf{Calibration Tab} featuring 4 synchronized subplots:
\begin{enumerate}
  \item \textbf{Raw Measurements:} Overlay of Pump (Air) black curve, Pump (Solvent) red curve, Single Mask cyan shaded area ($\alpha=0.4$), and Multiple Mask blue curve.
  \item \textbf{Calculated Absorbance \& Physical Baseline:} Purple absorbance spectrum with dashed orange background polynomial baseline.
  \item \textbf{Reference Overlay:} Reference FTIR spectrum overlaid against experimental absorbance, with zero-levels aligned via \code{mpl\_axes\_aligner}.
  \item \textbf{Calibration Fit:} Calibrated Wavelength (nm) vs.\ Pixel index dispersion curve with linear reference line and central crosshairs.
\end{enumerate}

The \textbf{Fit Controls} panel provides a global \emph{``Do baseline correction''} checkbox and defaults the \emph{``Fit Axis Unit''} selector to \textbf{Wavelength (nm)}. A standalone interactive Gaussian pump spectrum popup window allows interactive zoom/pan toolbar navigation and legend dragging.

\subsection{Programmatic Usage Example}

\begin{lstlisting}[language=Python]
import pymorgan.cal as cal

# 1. Load reference FTIR spectrum
ref = cal.load_reference_spectrum("FTIR-Dioxane.csv")

# 2. Load experimental probe spectrum (e.g. UniGE TRIR setup 9)
exp = cal.load_experimental_spectrum("pump_air.csv", cal_type_code=9)

# 3. Perform wavelength calibration fit
res = cal.fit_wavelength_axis(
    meas_y=exp.detector_data[0],
    ref_x=ref.spectral_axis,
    ref_y=ref.absorbance,
    cal_type_code=9,
    cwl=2000.0,
    use_derivative="1st",
)

print(f"Calibrated span: {res.wavelength_nm[0]:.2f} nm to {res.wavelength_nm[-1]:.2f} nm")

# 4. Save calibrated vector
cal.save_calibration_file(res.wavenumber_cm1, "output_dir", filename="CalibratedProbe.csv")

# 5. Fit probe beam Gaussian profiles (multi-detector)
probe_fit = cal.fit_probe_spectrum(exp.detector_data, exp.cwl)
for det_idx, fit in enumerate(probe_fit.fits):
    print(f"Detector {det_idx}: w0={fit.w0:.1f} cm-1, FWHM={fit.fwhm:.1f} cm-1")

# 6. Query default parameters for a given setup
params = cal.get_default_calibration_params(cal_type_code=9)
print(params["degree"])       # polynomial degree
print(params["cwl_guess"])    # initial central-wavelength guess
\end{lstlisting}

\section{Extending PyMORGAN (optional)}\label{sec:extending}

This optional chapter explains how to (i)~modify the GUI layout and behaviour
and (ii)~add a reader for a new instrument file format. Neither is required for
day-to-day use.

\subsection{Editing the GUI}
The window layout lives in a single Qt Designer file,
\code{pymorgan/gui/main\_window.ui}; the controller
\code{pymorgan/gui/main\_window.py} binds to the widgets \emph{by their
object name}. Editing the layout therefore means editing the \code{.ui} file,
and only \emph{new} widgets need code.

\subsubsection{Opening the layout}
Qt Designer edits the \code{.ui} regardless of the Qt binding. The easiest way to launch it is using the provided command:

\begin{lstlisting}[language=bash, style=py]
pymorgan-edit-gui
\end{lstlisting}

Alternatively, you can launch it manually using \code{uvx}:

\begin{lstlisting}[language=bash, style=py]
uvx --from pyside6-essentials pyside6-designer src/pymorgan/gui/main_window.ui
\end{lstlisting}

\subsubsection{Moving or resizing widgets}
Drag, resize and regroup freely. Because the controller looks widgets up by
object name (\code{getattr(self, "PP\_CtrPlot\_btn", \dots)}), rearranging the
layout breaks nothing as long as the object names are unchanged. Save the
\code{.ui} and re-run \code{pymorgan-gui}; no reinstall is needed because the
editable install reads the file at run time.

\subsubsection{Adding a widget}
\begin{enumerate}
  \item In Designer, add the widget and give it a clear \textbf{object name}
        (e.g.\ \code{PP\_exportFig\_btn}). After loading it is available as
        \code{self.PP\_exportFig\_btn}.
  \item In \code{main\_window.py}, inside \code{\_wire()}, connect it with the
        defensive helpers and write the slot:
\begin{lstlisting}
self._connect("PP_exportFig_btn", self.export_figure)   # buttons (clicked)
self._connect_text("SomeField", self.on_text_changed)   # QLineEdit (editingFinished)
\end{lstlisting}
        \code{\_connect}/\code{\_connect\_text} only attach if the widget
        exists, so adding-without-wiring (or removing a wired widget) never
        crashes the window.
  \item Read the widget's state in the slot (\code{self.<name>.currentText()},
        \code{.value()}, \code{.isChecked()}, \dots).
\end{enumerate}

\subsubsection{Plot areas and icons}
\begin{itemize}
  \item \textbf{A new plot area:} promote a \code{QWidget} to the class
        \code{WidgetPlot} with header \code{pymorgan.gui.widgetplot}
        (right-click~$\to$~\emph{Promote to\dots}), exactly like the existing
        \code{PPaxes}. In code it then exposes \code{.figure}, \code{.ax} and
        \code{.canvas}, so any pipeline plotter can draw into it via
        \code{data.plot\_contour(ax=widget.ax)}.
  \item \textbf{Icons:} do \emph{not} assign icons in Designer --- that embeds a
        compiled \code{.qrc} resource dependency that breaks \code{loadUi}.
        Instead drop the image in \code{pymorgan/gui/icons/} and add an entry to
        \code{\_BUTTON\_ICONS} in \code{main\_window.py} (applied by
        \code{\_apply\_icons}).
  \item Keep signal/slot wiring in Python (in \code{\_wire}), not in Designer's
        connection editor; the handlers live in \code{main\_window.py}.
\end{itemize}

\subsubsection{Do not rename}
The controller references these object names; renaming them in Designer requires
updating the matching strings in \code{main\_window.py}: \code{PPaxes},
\code{MainTabs}, \code{PP\_datatype\_cbx}, \code{PP\_datafolderlist\_lst},
\code{RootDir\_field}, the \code{PP\_*Plot\_btn} buttons, \code{PP\_SubtactBkg\_chk},
\code{PP\_bkg\_tmin}/\code{tmax} and \code{PP\_Normalise\_chk}. After any change,
confirm the window still builds:

\begin{lstlisting}[language=bash, style=py]
QT_QPA_PLATFORM=offscreen uv run pytest tests/test_gui.py
\end{lstlisting}

\subsection{Adding a new data-format parser}
A \emph{loader} maps a file path to the fields a dataset needs. Registering it
with the matching decorator makes its data-type name appear in
\code{pm.available\_*\_loaders()} and in the GUI's data-type combo.
Table~\ref{tab:loadercontracts} lists the three contracts.

\begin{table}[h]
  \centering
  \begin{tabular}{lll}
    \toprule
    Branch & Decorator & Returns \\
    \midrule
    1-D          & \code{@pm.register\_loader} & \code{(Zavg\_R, delays, probe, Units, Nscans, Zss\_R, Zstdv)} \\
    2-D          & \code{@pm.register\_map\_loader} & \code{(Z, pump, probe, delays, units, freq\_units)} \\
    steady-state & \code{@pm.register\_spectrum\_loader} & \code{(x, y, x\_units)} \\
    \bottomrule
  \end{tabular}
  \caption{Loader return contracts (see chapters~\ref{sec:oneD},~\ref{sec:twoD}).}
  \label{tab:loadercontracts}
\end{table}

\subsubsection{A 1-D reader}
The signal is stored as \code{[Ndelays $\times$ Npixels $\times$ Ndetectors]};
the unit dictionary is built with \code{helpers.units2dic(unitsL, unitsT, unitsZ)}.
A minimal comma-separated reader:

\begin{lstlisting}
import numpy as np
import helpers as hlp
import pymorgan as pm

@pm.register_loader("MyTA")
def read_my_ta(path):
    raw = np.loadtxt(path, delimiter=",")
    delays = raw[1:, 0]                       # first column  = delays
    probe  = raw[0, 1:]                        # first row     = probe axis
    Zavg   = raw[1:, 1:][..., np.newaxis]      # add a detector axis
    Units  = hlp.units2dic("nm", "ps", "mOD")  # (probe, time, signal) units
    Nscans = np.nan                            # no single-scan information
    Zss    = np.nan * np.zeros_like(Zavg[..., np.newaxis])
    Zstdv  = np.zeros_like(Zavg)
    return Zavg, delays, probe, Units, Nscans, Zss, Zstdv
\end{lstlisting}

The data type is then available everywhere:

\begin{lstlisting}
data = pm.load_1D("scan.myta", data_type="MyTA")
\end{lstlisting}

\paragraph{Where to register.}
The decorator runs when its module is imported. For a permanent format, place the
reader in \code{pymorgan/oneD/load.py} (imported with the package). For a one-off,
defining it anywhere before calling \code{pm.load\_1D} is enough.

\paragraph{Exposing it in the GUI.}
Add the file extension to \code{\_DATATYPE\_GLOB} in
\code{pymorgan/gui/main\_window.py} so the dataset browser lists matching files,
for example \code{"MyTA": "*.myta"} (unlisted types fall back to \code{"*"}).

\subsubsection{2-D and steady-state readers}
The pattern is identical, only the returned tuple differs
(Table~\ref{tab:loadercontracts}). A 2-D reader returns the cube
\code{Z[pump, probe, t2]} with its pump/probe/t2 axes, a unit dictionary and the
spectral-axis unit string; a steady-state reader returns the \code{x}/\code{y}
arrays and a best-guess spectral unit (or \code{None}). Register them with
\code{@pm.register\_map\_loader("Name")} and
\code{@pm.register\_spectrum\_loader("Name")} respectively.

\subsubsection{The unit dictionary}
\code{helpers.units2dic(unitsL, unitsT, unitsZ)} accepts the spectral unit
\code{unitsL} (\code{"nm"}, \code{"cm\textasciicircum\{-1\}"} or \code{"eV"}),
the time unit \code{unitsT} (e.g.\ \code{"ps"}), and the signal unit
\code{unitsZ} (\code{"mOD"}, \code{"x1E3"} or \code{"int"} for intensity). The
\dA{} display convention (m\,OD vs $\times 10^{3}$) is then a presentation choice
handled by the settings (\S\ref{sec:settings}).


\section{License}\label{sec:license}
PyMORGAN is free software released under the \textbf{GNU Affero General Public License version~3.0} (\textsf{AGPLv3}).
You are free to run, inspect, modify and redistribute this software under the terms of the license. Any derivative works or modifications made available over a network or server must also make their complete corresponding source code available under the same license terms.

\end{document}
